\subsection{Proof of Theorem~\ref{thm:ind_noise}}
The coefficient condition implies $.999\le c(v)\le1.002$ for every $v\ge v_L$.
We use the margins in Section~\ref{sec:geometry_full} with the separate coordinate tolerances
in~\eqref{eq:ind_event}. At the barrier's right face, the horizontal
tolerance increases by at most $.001/4$, so $q_B>.298896-.00025/.999>.29864$.
The certified bounds $L(.17)<.287729$ and $U(2.5)<.289884$ still suffice.
The bottom barrier and channel tests have left side at least
$\gamma[(1+.999)(.01)-.005]>10^{-4}$.
For channel progress, $b_1-\alpha s\ge \alpha(1-s_E-1/4)>0$.
The comparison sequence uses this sign and the same $\kappa,\Delta_E$,
and so the 25-step bound is unchanged. The initial $s$ is between
$s_-=-.01$ and $s_E=.53$ for class 1, and $|t|\le.005$.
The certified channel top-face drift divided by $\gamma$ increases
by at most $.005$, from at most $-.01064024$.
The overshoot bound increases by at most $\gamma(.00025)$, remaining
below $.649<s_U$. The right-face test holds since
$\gamma[(\alpha+.999)(.1)-.00025]>10^{-4}$.
The classifier thresholds still distinguish the three means.
Class 0 starts within $.005$ of the origin in each coordinate, and
the class-2 initial coordinates satisfy $-.00127\le s\le .00025$
and $1.975\le t\le2.025$, inside $\widetilde G_2$.

The terminal gaps remain at least $.06$ and imply
$\Lambda_z\le4e^{-.03/v_L}\le4e^{-30}$.
Every template-to-face distance is at least $.09$, so
$.005+1.002\cdot3\cdot4e^{-30}+10^{-4}/\gamma<.999\cdot .09$.
The terminal invariance test holds with the global error budget.
In the local residual bound the linear factor is at most $1-\gamma(.999)<q$.
The entry distance to $m_Z(c_L)$ is at most 6, and
$c_L/(\alpha+c_L)\le1$. Hence on $E$ the residual norm after at least $n$
terminal calls is at most
\[
 6q^n+12e^{-.03/v_L}+(\sqrt{2}/(\gamma\cdot.999))\delta_L .
\]

Conditional on $Z$, $e=Bu+A^\top\eta$ has coordinate variances
$\alpha^2\tau^2+\alpha\nu^2$ and $\tau^2+\nu^2$ and its law is independent of $Z$.
Gaussian tails and the union bound give~\eqref{eq:ind_tail}.
On $E$ for $Z=1$, the restricted output satisfies $1-\widehat s\ge .7$.
Applying the reverse triangle inequality in conditional $L^2$ to
$(1-\widehat s+u_1)\mathbf{1}_E$ gives~\eqref{eq:ind_lower}.

For the upper bound, the exact coordinate update is a convex combination
of $x_j$ and $(b_j+c(v)\bar\mu_{v,j}(x))/(\beta_j+c(v))$, with weight
$\gamma(\beta_j+c(v))\ge .999/4$ on the latter.
The latter value has absolute value bounded by $3+2|e_j|$, using
$b_j=\beta_j\mu_{Zj}+e_j$ and the bounded template coordinates.
The same bound holds at initialization and for $m_Z(c_L)$.
With update errors, induction adds at most $10^{-4}/(\gamma\cdot.999)$
to the bound on the iterates. As $\operatorname{Var}(e_j)\le S^2\le1$,
Gaussian fourth moments give a uniform $L^4$ bound on
$\|\widehat x-m_Z(c_L)\|$, independent of $\alpha$ and of the budget.
By Cauchy--Schwarz the contribution to the squared residual on $E^c$
is at most a constant times $p_E^{1/2}$.
Finally,
\[
 \E\norm{m_Z(c_L)-x^\star}^2
   =\sum_{\beta\in\{\alpha,1\}}
     \frac{\beta\nu^2+c_L^2\tau^2}{(\beta+c_L)^2}
   \le C_0 S^2 .
\]
The squared triangle inequality gives~\eqref{eq:ind_upper}.
Under~\eqref{eq:ind_regime}, $p_E\le4S^4$, the terminal leakage is at most
$S$ before squaring, and $q^n\le S$. Since $\pi_1>.30$, the stated
separation follows. At fixed $\alpha>0$, both Gaussian tail conditions
in~\eqref{eq:ind_regime} hold for sufficiently small $S$. Also $\tau^2/v_L\le S$ for $v_L=S$, which yields the coefficient window
eventually at $\lambda=1$.

